% ===== §23 =====
\section{The Special Configuration of the Three Registers}
\label{p24:sec:registers}

\noindent
The fourth pillar returns the three registers of the framework to the intimate dyad and shows that here, and nowhere else among cultural formations, they attain their fullest and most intense configuration. The meta-theoretical statement claimed that culture spans the imaginary, the symbolic, and the real, and that the three move in a materially driven dialectic; the intimate relation is the site where this claim is most completely embodied, because in no other cultural formation are the three registers configured with the closeness, the intensity, and the mutual pressure they have between two people who share a life. This pillar describes that configuration register by register, and then shows the dialectic among them running at a pitch the macro-scale never reaches.

\subsection*{The three registers at their most intense}

The \emph{imaginary} is, in intimate relation, unusually thick. The register of image and identification, of the body and the specular relation, is nowhere more charged than between two people bound by desire: each is an image for the other of extraordinary power, each sees themselves reflected and altered in the other's regard, and the bodily, visual, identificatory life of the couple---how they appear to each other, how each is formed by the other's image of them---is intense in a way no macro-culture's imaginary approaches. A nation's imaginary, its images and identifications, is diffuse and mediated; a couple's is concentrated in a single face.

The \emph{symbolic} is, as the earlier pillars showed, the register of the private lexicon, the shared references, the names and rites and promises through which the interwoven desire is given sayable form. Here the intimate symbolic reaches its highest instance in the vow---the promise two people make to each other, which this series has treated at length as the non-binding vow whose whole force is that it appeals to no external enforcement but institutes, between the two, a symbolic order they legislate themselves \parencite{paperxiii}. The couple's symbolic is not a fragment of the public symbolic order but a private one, coined and instituted in the vacancy of the big Other, and the vow is its constitutional act.

The \emph{real} is, in intimate relation, exposed as it is almost nowhere else. The register of the unsymbolizable is present in the couple's life in its most naked forms: in sexuality, where the drive meets what no signifier orders; in the traumatic, where the bond meets what cannot be assimilated; and above all in the ``us'' itself, the non-symbolizable core of the bond that no private word quite captures, which is the relational Real this paper has placed at the center of its ontology. Intimate relation is one of the few sites of human life where the Real is not held at a distance by a thick symbolic mediation but is met closely, in the body and in the bond, so that the relational Real is not a theoretical postulate here but a lived and constant presence.

This grounds the fifth thesis the paper advances as its own, which gathers into a single claim what the earlier theses have each drawn on. Call it the \emph{real-anchoring thesis}: the resistance a relational culture offers to power, to capital, and to time arises from its anchoring in the register of the Real, and the firmness of that anchoring rises with the thickness of what, in the participants' interwoven drive, exceeds symbolization. Resistance, on this thesis, is not an added property a culture might have or lack, a defensive feature bolted onto its forms; it is the persistence of the culture's root in the Real under the symbolic erosion that power and capital work. Where a culture is anchored only in the symbolic---in convention, in transferable form---it offers little resistance, because the symbolic is precisely the register in which power organizes and capital extracts; where it is anchored in an interwoven Real that exceeds symbolization, it resists, not by opposing a force to these operations but by lying, in part, outside the register on which they operate. The thesis does not equate resistance with anchoring depth as one quantity with another; it locates the source of the one in the other, and makes the degree of resistance vary with the thickness of the unsymbolizable rather than fixing it. The other four theses can be read through this one: non-transferability, decoherence, the vacancy's fragility, and the counter-power of persistence all turn on where, and how deeply, a culture is anchored in a Real that the symbolic operations of transfer, capture, guarantee, and entrenchment cannot reach.

\subsection*{The dialectic of the registers, at full intensity}

Because the three registers are configured so closely and so intensely in the dyad, the dialectic among them---which the fourth revision set at the center of the framework---runs here at a pitch found nowhere else, and it can be watched in the couple's culture with unusual clarity. The dialectic, in its direct form, is the contest between the symbolic's drive to capture the real and the real's overflow of every symbolic capture. In intimate relation this contest is continuous and intense. Each private word, each name, each rite, each vow is an attempt by the couple's symbolic to capture the ``us,'' to give the unsymbolizable bond a sayable and repeatable form---and each such capture is, necessarily, incomplete, because the relational Real of the bond overflows every form the couple can give it. The lovers name the bond and the name is precious and the name is not the bond; they institute a rite to hold what passes between them and the rite holds some of it and not all; they make a vow that says what they are to each other and the vow reaches toward but does not exhaust the real of their being-for-each-other. This incompleteness is not a failure to be corrected by a better word; it is the structural condition the fourth revision named, the Real's constitutive overflow of the symbolic, and in intimate relation it is felt directly, as the perpetual insufficiency of every expression to the bond it would express, and the perpetual generation of new expressions to chase what the last one could not hold. The couple's culture is never finished precisely because the relational Real it tries to capture overflows it, and this is why intimate culture develops, continually, generating new forms as the old ones prove unequal to the ``us'' they were made to hold.

\subsection*{Why the dyad embodies the framework most fully}

This is the sense in which the intimate dyad is the fullest embodiment of the framework's central claim, and marking it completes the pillar. The three-register dialectic of the relational Real was stated in Part One as a general ontology of culture, holding of culture at any scale. But it holds of the intimate dyad with an intensity and a directness it has nowhere else, because in the dyad the three registers are configured at their maximum---the imaginary concentrated in a single face, the symbolic legislated privately in the vacancy of the big Other, the real met nakedly in body and bond---and the dialectic among them, the symbolic's capture and the real's overflow, runs at a pitch that makes it visible in the ordinary life of the couple. What in macro-culture is a slow and buried process is, in the dyad, an intense and nearly witnessable one: the perpetual chase of the symbolic after a relational Real that will not be caught, generating, in its perpetual insufficiency, the perpetual development of the couple's culture. The framework was built to be general; the dyad is where it is most fully true. And the same overflow that keeps the culture developing---the Real's refusal to be captured---is what will, in the register of loss, expose the culture to the collapse the final pillar must now describe.
